\maketitle

\begin{abstract}
A $\sigma$-additive probability measure on an observable
algebra~\cite{mahon_paper1} already determines, via the Rokhlin
disintegration theorem, a disintegration kernel for any pair of
observables and a minimal sufficient factor carrying full conditional
information.  When the observables are indexed over time, temporal
coherence forces the Chapman--Kolmogorov equation; a Markov semigroup
and Koopman--Perron duality follow without a separately posited
dynamical law.  In a measure-preserving system the minimal sufficient
factor becomes the delay map, and reconstruction --- that the observable
$\sigma$-algebra equals the full $\sigma$-algebra modulo null sets, that
the delay algebra is dense in $L^2$, and that the delay map is a
measure-theoretic embedding --- reduces to the faithfulness of this
factor.  The Stone space built in~\cite{mahon_paper1} to realise
$\sigma$-additivity is then identified with the state space itself.
\end{abstract}

\begin{quote}\small
\textbf{2020 Mathematics Subject Classification.}
37A30 (primary); 47D06, 60J05, 37A05, 28A60, 06E15, 46E30 (secondary).

\textbf{Keywords.}
Koopman operator, disintegration kernel, minimal sufficient factor,
Markov semigroup, Koopman--Perron duality, observable dynamical system,
reconstruction theorem, faithful factor, delay embedding, observable algebra,
density bridge, generating partition, Stone duality, measure-preserving system.
\end{quote}

% ====================================================================
\section{Introduction}
\label{II:sec:intro}

Koopman studied dynamical systems through observables~\cite{Koopman1931}, but
the dynamical law $T$ remains primitive: the operator $U_T f = f \circ T$ is
defined only after an underlying evolution has already been specified.  The
classical stochastic-process tradition proceeds similarly.  Prediction is
formulated relative to a filtration, so temporal order and information flow are
built into the formalism from the outset~\cite{Doob1953}.  That ordering is
natural, but it need not be conceptually prior.  One may instead ask whether observable probability already carries a conditional
regularity between observables, and whether the familiar semigroup and predictive
structures arise only when that regularity is organised coherently over time.

In~\cite{mahon_paper1}, $\sigma$-additivity is the irreducible commitment
that turns coherent charges into a probability measure: no first-order
condition can force it, and compatibility collapses every candidate
reformulation back to $\sigma$-additivity itself.  The Stone construction
produces a measure unconditionally, but only $\sigma$-additivity ensures it
lives on the space of realised states.  Once committed to, however, the
measure already carries more than it was built to carry.  Any two
observations --- no hypothesis on time or order --- determine a
conditional regularity: the Rokhlin disintegration
theorem~\cite{Rokhlin1952} forces a Markov kernel into existence,
from which a minimal sufficient factor emerges as the coarsest
reduction of the sample point carrying full conditional information.

Index this regularity over time and the kernels must compose
consistently.  That coherence forces the Chapman--Kolmogorov
equation, yielding a Markov semigroup and Koopman--Perron duality
without a separately posited dynamical law.  In a measure-preserving
system the minimal sufficient factor becomes the delay map, and the
question becomes internal: does the observable algebra generated by
the delay orbit equal the full $\sigma$-algebra?  The reconstruction
theorem gives three equivalent formulations; under any of them, the
Stone space~\cite{mahon_paper1} built to realise $\sigma$-additivity
collapses onto the state space itself.

% ====================================================================
\section{Disintegration, dynamics, and duality}
\label{II:sec:prediction}


Let $(\Omega, \mathcal{F}, P)$ be a $\sigma$-additive probability space on an
observable algebra~\cite{mahon_paper1}, and let $Q, F$ be two measurable maps:
\begin{itemize}
  \item $Q : \Omega \to \alpha$, where $(\alpha, \mathcal{A})$ is a measurable space;
  \item $F : \Omega \to \beta$, where $(\beta, \mathcal{B})$ is a standard Borel space.
\end{itemize}
The standard Borel hypothesis on $\beta$ is used once: it is the hypothesis under
which the Rokhlin disintegration theorem~\cite{Rokhlin1952} guarantees the existence
of a regular conditional distribution.  All other results hold in greater generality.

We write $\mathcal{P}(\beta)$ for the space of probability measures on $\beta$,
equipped with its canonical (Giry) $\sigma$-algebra.  The joint pushforward
$\rho := P \circ (Q, F)^{-1}$ on $\alpha \times \beta$ has marginals
$\rho_\alpha = P \circ Q^{-1}$ and $\rho_\beta = P \circ F^{-1}$.

The regular conditional distribution of $F$ given $Q$ is the unique
(up to $P \circ Q^{-1}$-null sets) Markov kernel
$\kappa_Q : \alpha \to \mathcal{P}(\beta)$ satisfying
\[
  \int_A \kappa_Q(q, B)\, d(P \circ Q^{-1})(q) \;=\; P(Q \in A,\; F \in B)
\]
for all measurable $A \subseteq \alpha$ and $B \subseteq \beta$.

\begin{definition}[Disintegration kernel]
\label{II:def:predictive-kernel}
The \emph{disintegration kernel} of $F$ given $Q$ is $\kappa_Q$, the
regular conditional distribution of $\rho = P \circ (Q, F)^{-1}$ over
its first marginal:
$\kappa_Q(q, A) = P(F \in A \mid Q = q)$.
\end{definition}

The kernel is covariant under measurable maps: if $\pi : \alpha \to \alpha'$ is
measurable and $Q' = \pi \circ Q$, then $\kappa_Q(q, \cdot) = \kappa_{Q'}(\pi(q), \cdot)$
for $(P \circ Q^{-1})$-almost every $q$, by uniqueness of disintegration.

\begin{definition}[Minimal sufficient factor]
\label{II:def:minimal-predictive}
The \emph{minimal sufficient factor} is
\[
  Q_* : \Omega \to \mathcal{P}(\beta), \qquad Q_*(\omega) \;=\; \kappa_Q(Q(\omega), \cdot).
\]
\end{definition}

Since $\kappa_Q$ is a Markov kernel it is measurable, and $Q_*$ inherits
measurability by composition.

\begin{theorem}[Factorization via the canonical factor]
\label{II:thm:predictive-factorization}
For every bounded measurable $g : \beta \to \mathbb{R}$:
\[
  E[g(F) \mid \sigma(Q)] \;=\; \int g\, dQ_*(\cdot)
  \quad P\text{-a.e.}
\]
\end{theorem}
\begin{proof}
Let $h(\omega) := \int g\, dQ_*(\omega)$.  We verify the two defining properties.

\emph{Measurability.}  The map $q \mapsto \int g\, d\kappa_Q(q, \cdot)$ is
strongly measurable; since $Q_* = \kappa_Q(Q(\cdot), \cdot)$, the map $h$ is
$\sigma(Q)$-measurable.

\emph{Set-integral identity.}  For $S = Q^{-1}(A)$:
\begin{align*}
  \int_S h\, dP
  &= \int_A \int g\, d\kappa_Q(q, \cdot)\, d(P \circ Q^{-1})(q)
  \quad\text{(pushforward)} \\
  &= \int_{A \times \beta} g(f)\, d\rho(q, f)
  \quad\text{(disintegration)} \\
  &= \int_S g(F)\, dP.
\end{align*}
The conclusion follows from a.e.\ uniqueness of conditional expectation.
\end{proof}

Since $\sigma(Q_*) \subseteq \sigma(Q)$ and the right-hand side is already
$\sigma(Q_*)$-measurable, $Q_*$ is a sufficient statistic for $F$ given $Q$:
$E[g(F) \mid \sigma(Q)] = E[g(F) \mid \sigma(Q_*)]$ $P$-a.e.
It is \emph{minimal} in the sense that every sufficient statistic factors through
it: $Q_*$ is the coarsest reduction of $\omega$ that carries the full conditional
regularity of $F$ given $Q$.
The operator $(K_{Q_*} g)(q_*) = \int g\, dq_*$ on $\mathcal{P}(\beta)$
recasts this as $E[g(F) \mid \sigma(Q)] = (K_{Q_*} g) \circ Q_*$ $P$-a.e.


When the disintegration kernel $\kappa_Q$ is indexed over a time parameter,
the kernels cannot be chosen independently: the regularity of $F$ at time $t+s$
given $Q$ at time $0$ must be consistent with first observing at time $s$ and
then at time $t$.  This consistency requirement forces the
Chapman--Kolmogorov equation.

Concretely, index over a discrete time horizon $\mathbb{N}$ and let
$\Pi_t : \gamma \to \mathcal{P}(\gamma)$ be a Markov kernel on a state
space $\gamma$, with $\Pi_0(q, \cdot) = \delta_q$.  Temporal coherence of
the conditional regularity demands
\[
  \Pi_{t+s}(q, \cdot) \;=\; \int_\gamma \Pi_t(q', \cdot)\, d\Pi_s(q, q')
  \quad\text{(i.e., } \Pi_{t+s} = \Pi_t \circ_\kappa \Pi_s\text{)}
\]
— the Chapman--Kolmogorov equation, derived from the structure rather than
assumed.  A family $\{\Pi_t\}$ satisfying it is a \emph{Markov semigroup kernel}.
The associated \emph{Markov operators} are
\[
  (K_t g)(q) \;:=\; \int g\, d\Pi_t(q, \cdot), \qquad g : \gamma \to \mathbb{R}.
\]

Each $K_t$ is positive, unital, and contractive on bounded measurable functions,
with $K_0 = \mathrm{id}$; positivity, unitality, and contractiveness are immediate
from the Markov property.

\begin{proposition}[Semigroup property]
\label{II:prop:semigroup}
$K_{t+s} g = K_t(K_s g)$ for all bounded measurable $g$.
\end{proposition}
\begin{proof}
$(K_{t+s} g)(q)
= \int g\, d(\Pi_t \circ_\kappa \Pi_s)(q, \cdot)
= \int \bigl(\int g\, d\Pi_t(q', \cdot)\bigr)\, d\Pi_s(q, q')
= \int (K_t g)(q')\, d\Pi_s(q, q')
= (K_t(K_s g))(q)$,
where the exchange of integrals is justified by Bochner--Fubini.
\end{proof}

The pushforward of a measure $\mu$ on $\gamma$ through the kernel is
$P_t^* \mu := \mu \star \Pi_t = \int \Pi_t(q, \cdot)\, d\mu(q)$.


\begin{theorem}[Koopman--Perron duality]
\label{II:thm:koopman-perron}
For every bounded measurable $g$ and finite measure $\mu$:
\[
  \int (K_t g)\, d\mu \;=\; \int g\, d(P_t^* \mu).
\]
\end{theorem}
\begin{proof}
$\int g\, d(P_t^* \mu)
= \int g\, d(\mu \star \Pi_t)
= \int_\gamma \int g\, d\Pi_t(q, \cdot)\, d\mu(q)
= \int (K_t g)\, d\mu$,
by Bochner--Fubini.
\end{proof}

When $\Pi_t(q, \cdot) = \delta_{\varphi_t(q)}$ for a measurable map
$\varphi_t : \gamma \to \gamma$, the stochastic picture collapses.


\begin{proposition}[Deterministic limit]
\label{II:prop:deterministic}
If $\Pi_t(q, \cdot) = \delta_{\varphi_t(q)}$, then $K_t g = g \circ \varphi_t$
and $\varphi_{t+s} = \varphi_t \circ \varphi_s$.
\end{proposition}
\begin{proof}
$(K_t g)(q) = \int g\, d\delta_{\varphi_t(q)} = g(\varphi_t(q))$.
The flow law follows from
$\delta_{\varphi_{t+s}(q)} = (\Pi_t \circ_\kappa \Pi_s)(q, \cdot)
= \delta_{\varphi_t(\varphi_s(q))}$
and injectivity of the Dirac embedding.
\end{proof}

In the deterministic case, $K_t g = g \circ \varphi_t = U_T^t g$: the Markov
operator specialises to the classical Koopman operator.  The
Koopman--Perron duality $\int K_t g\, d\mu = \int g\, d(P_t^* \mu)$ becomes
the change-of-variables formula $\int g \circ T^t\, d\mu = \int g\, d(T^t_* \mu)$.

\begin{example}[Finite Markov chain]
\label{II:ex:markov}
Let $X = \{1,\ldots,k\}$, $\mathcal{B} = 2^X$, and let $P = (p_{ij})$ be the
transition matrix of a finite irreducible Markov chain with unique stationary
distribution $\mu$ ($\mu P = \mu$).  Take full-state observation $Q_t = X_t$.
The disintegration kernel of $X_1$ given $X_0$ is
$\kappa(i,\cdot) = P(i,\cdot)$, the $i$-th row of $P$: this is the unique Markov
kernel satisfying
$\mathbb{E}_\mu[f(X_1) \mid X_0 = i] = \sum_j p_{ij} f(j)$ for all bounded $f$.
The $t$-step kernel $\Pi_t(i,\cdot) = P^t(i,\cdot)$ is the corresponding
semigroup kernel, and temporal coherence gives
$\Pi_{t+s} = \Pi_t \circ_\kappa \Pi_s$, i.e.\ $P^{t+s} = P^t P^s$.
The semigroup acts as
\[
  K_t f(i) = \sum_j p^t_{ij} f(j).
\]
Koopman--Perron duality reduces to the stationary-measure identity:
\[
  \int K_t g\, d\mu = \mu^\top P^t g = (\mu P^t)^\top g = \mu^\top g
  = \int g\, d\mu,
\]
the fixed-point property $\mu P = \mu$ of the stationary distribution.
In this finite setting every charge is automatically $\sigma$-additive, so the
admissibility condition of~\cite{mahon_paper1} is trivially satisfied.
Reconstruction holds at lag $0$: $\mathcal{O}_h = \sigma(h) = \mathcal{B}(X)$
whenever $h$ separates points of $X$.
\end{example}

The full structure produced by the construction is the triple
$(\mathcal{P}(\beta),\, \nu_{Q_*},\, \{K_t\}_{t \in \mathbb{N}})$:
the canonical factor space, stationary distribution $\nu_{Q_*} = P \circ Q_*^{-1}$,
and Markov semigroup.  This is the \emph{observable dynamical system}
derived from $(Q, F, P)$: the conditional regularity of $F$ relative to $Q$,
promoted to a dynamical object by temporal coherence.

% ====================================================================
\section{Reconstruction}
\label{II:sec:reconstruction}

The structures of §\ref{II:sec:prediction} require only a probability
measure and two observables.  We now specialise: let
$(X, \mathcal{B}, \mu, T)$ be an invertible measure-preserving system and
$h \in L^\infty$ an observable.  Setting $Q_t = h \circ T^t$ gives
$Q_*(\omega) = (h(T^n \omega))_{n \in \mathbb{Z}}$, the delay map
$\Phi_h(x) = (h(T^n x))_{n \in \mathbb{Z}}$, and the observable
dynamical system $(\Phi_h(X), \mu \circ \Phi_h^{-1}, \{K_t\})$.
The measure-preserving hypothesis and the choice of $h$ are genuine
inputs; what follows from them is not.

The canonical factor space $\Phi_h(X)$ is a subset of $\mathbb{R}^\mathbb{Z}$
encoding all dynamical information accessible to $h$.  Whether it is a faithful
image of $X$ depends on whether the delay orbit
$\{h \circ T^n : n \in \mathbb{Z}\}$ separates points of $X$.
Injectivity of $\Phi_h$ modulo $\mu$-null sets is equivalent to the observable
$\sigma$-algebra
\[
  \mathcal{O}_h = \sigma\!\bigl(\{h \circ T^n : n \in \mathbb{Z}\}\bigr)
\]
equalling $\mathcal{B}(X)$ modulo $\mu$.

\begin{definition}[Faithful factor]
\label{II:def:faithful}
A measurable map $\psi : (X, \mathcal{B}, \mu) \to (Y, \mathcal{Y})$ is
\emph{faithful modulo $\mu$} if
\[
  \sigma(\psi) \;=\; \mathcal{B} \quad \text{modulo } \mu,
\]
equivalently, if $\psi$ separates $\mu$-a.e.\ pair of points of $X$.
The reconstruction question is whether $\Phi_h$ is faithful modulo~$\mu$.
\end{definition}

The Koopman--Perron duality (Theorem~\ref{II:thm:koopman-perron}) connects this to a
function-space question: the observable $\sigma$-algebra $\mathcal{O}_h$ determines
the Markov semigroup through the algebra it generates, and the reconstruction question
is precisely whether that algebra is large enough to recover all of $\mathcal{B}$.


For a finite joined query algebra $\mathcal{G} = \mathcal{E}_{n_1} \vee \cdots
\vee \mathcal{E}_{n_k}$ from a directed system $(\mathcal{E}_n)_n$, the
\emph{separation defect} is
\[
  \delta(\mathcal{G})
  \;:=\;
  \sup_{A \in \sigma(\mathcal{C})} \inf_{E \in \mathcal{G}} \mu(A \triangle E),
\]
where $\sigma(\mathcal{C})$ is the observable $\sigma$-algebra of the full
system.  It vanishes iff $\sigma(\mathcal{G}) = \sigma(\mathcal{C})$ mod $\mu$.
A sequence $(\mathcal{G}_k)$ is \emph{honest} if it is increasing and
$\sigma(\bigcup_k \mathcal{G}_k) = \sigma(\mathcal{C})$ mod $\mu$; under an
honest scheme $\delta(\mathcal{G}_k)$ is non-increasing.

\begin{theorem}[CE drives separation defect to zero]
\label{II:thm:ce-sep-defect}
Let $(\mathcal{E}_n)_n$ be a query system with realization measure $\mu$
satisfying collective exhaustion~\cite{mahon_paper1}.  Under any honest
refinement scheme $(\mathcal{G}_k)$,
\[
  \delta(\mathcal{G}_k) \;\longrightarrow\; 0
  \quad \text{as } k \to \infty.
\]
\end{theorem}

\begin{proof}
Fix $A \in \sigma(\mathcal{C})$.  Since $(\mathcal{G}_k)$ is honest,
$\sigma(\bigcup_k \mathcal{G}_k) = \sigma(\mathcal{C})$ mod $\mu$, so $A$
is $\sigma(\bigcup_k \mathcal{G}_k)$-measurable.  By the $L^1(\mu)$
martingale convergence theorem, $\mathbf{1}_{E_k} \to \mathbf{1}_A$ in
$L^1(\mu)$ where $E_k = \mathbb{E}_\mu[\mathbf{1}_A \mid \mathcal{G}_k]$
rounded to the nearest atom; in particular $\mu(A \triangle E_k) \to 0$.
CE~\cite{mahon_paper1} ensures no mass persists on the vanishing sets
$A \triangle E_k$, confirming $\inf_{E \in \mathcal{G}_k}\mu(A \triangle E)
\to 0$ uniformly over $A$.  Hence $\delta(\mathcal{G}_k) \to 0$.
\end{proof}

\begin{remark}[The temporal specialisation]
\label{II:rem:temporal-spec}
With $Q_t = h \circ T^t$, the algebra $\mathcal{O}_h^{(L)}$ is an honest
refinement scheme once reconstruction holds.
The separation defect specialises to $\delta(L) := \delta(\mathcal{O}_h^{(L)})$;
under mixing, $\delta(L)$ decays exponentially.  In the general
query setting, exhaustion and mixing are separate conditions; the delay-map
setting collapses the first by construction.
\end{remark}


Throughout, $(X, \mathcal{B}, \mu)$ is a standard
probability space, $T : X \to X$ is a measurable, $\mu$-preserving invertible
bimeasurable bijection with $T_*\mu = \mu$, and $h \in L^\infty(X, \mu)$.

\begin{definition}[Observable algebra]
\label{II:def:obs-algebra}
The \emph{observable algebra} generated by $h$ is
\[
  \mathcal{O}_h \;:=\; \sigma\!\bigl(\{h \circ T^n : n \in \mathbb{Z}\}\bigr)
  \;\subseteq\; \mathcal{B}(X).
\]
The \emph{delay algebra} $\mathcal{A}_h$ is the smallest subalgebra of
$L^\infty(X, \mu)$ containing $1$ and every $h \circ T^n$; concretely, finite
sums $\sum_\alpha c_\alpha \prod_k (h \circ T^{n_k})^{e_k}$.
\end{definition}

\begin{definition}[Delay map]
\label{II:def:delay-map}
The \emph{delay map} is
\[
  \Phi_h : X \to \mathbb{R}^\mathbb{Z},
  \qquad
  \Phi_h(x) \;=\; \bigl(h(T^n x)\bigr)_{n \in \mathbb{Z}}.
\]
Since $\pi_n \circ \Phi_h = h \circ T^n$ for each coordinate projection $\pi_n$,
the map $\Phi_h$ is measurable and
$\mathcal{O}_h = \Phi_h^{-1}(\mathcal{B}(\mathbb{R}^\mathbb{Z}))$.
Thus $\mathcal{O}_h = \mathcal{B}$ mod $\mu$ is precisely the condition that
$\Phi_h$ is a measure-theoretic embedding.
\end{definition}

The central lemma connects the Boolean algebra level to the $L^2$ level.
It holds for any subalgebra of $L^\infty$.

\begin{lemma}[Density bridge]
\label{II:lem:density-bridge}
Let $(X, \mathcal{B}, \mu)$ be a probability space and let
$\mathcal{A} \subseteq L^\infty(X, \mu)$ be a subalgebra containing the
constant function~$1$.  Then:
\[
  \overline{\mathcal{A}}^{L^2}
  \;=\;
  L^2(X,\, \sigma(\mathcal{A}),\, \mu).
\]
Consequently, $\overline{\mathcal{A}}^{L^2} = L^2(X, \mu)$ if and only if
$\sigma(\mathcal{A}) = \mathcal{B}$ modulo $\mu$.
\end{lemma}

\begin{proof}
Write $\mathcal{M} := \sigma(\mathcal{A})$ and $V := \overline{\mathcal{A}}^{L^2}$.

\textbf{Step 1: $V \subseteq L^2(X, \mathcal{M}, \mu)$.}
Every $f \in \mathcal{A}$ is $\mathcal{A}$-measurable, hence $\mathcal{M}$-measurable.
So $\mathcal{A} \subseteq L^2(X, \mathcal{M}, \mu)$.  Since $L^2(X, \mathcal{M}, \mu)$
is a closed subspace of $L^2(X, \mathcal{B}, \mu)$, taking the closure gives
$V \subseteq L^2(X, \mathcal{M}, \mu)$.

\textbf{Step 2: $L^2(X, \mathcal{M}, \mu) \subseteq V$.}
It suffices to show that every indicator $\mathbf{1}_E$ with $E \in \mathcal{M}$
lies in $V$.  By the functional monotone class theorem~\cite{Dynkin}, the class
\[
  \mathcal{H}
  \;:=\;
  \bigl\{f \in L^\infty(X, \mathcal{M}, \mu) : f \in V\bigr\}
\]
is a vector space containing $\mathcal{A}$ and closed under bounded monotone
pointwise limits (since if $f_n \in V$ are uniformly bounded and $f_n \to f$
pointwise a.e.\ with $f \in L^\infty$, then $f_n \to f$ in $L^2$ by dominated
convergence, so $f \in V$).  Since $\mathcal{A}$ is an algebra containing $1$
and $\sigma(\mathcal{A}) = \mathcal{M}$, the monotone class theorem gives
$\mathcal{H} = L^\infty(X, \mathcal{M}, \mu)$.  In particular every
$\mathcal{M}$-indicator $\mathbf{1}_E \in V$.

Since the $\mathcal{M}$-simple functions are dense in $L^2(X, \mathcal{M}, \mu)$
and $V$ is closed, $L^2(X, \mathcal{M}, \mu) \subseteq V$.

\textbf{Conclusion.}
Steps~1 and~2 give $V = L^2(X, \mathcal{M}, \mu)$.  The final equivalence is
immediate: $V = L^2(X, \mathcal{B}, \mu)$ iff $L^2(X, \mathcal{M}, \mu) =
L^2(X, \mathcal{B}, \mu)$ iff $\mathcal{M} = \mathcal{B}$ mod $\mu$.
\end{proof}

\begin{remark}[Role of the algebra structure]
\label{II:rem:algebra-necessary}
The algebra structure of $\mathcal{A}$ is essential.  The corresponding
statement for the closed \emph{linear} span $\mathcal{H}_{\mathcal{A}} =
\overline{\mathrm{span}}(\mathcal{A})$ is false in general:
$\mathcal{H}_{\mathcal{A}} = L^2$ does not imply
$\sigma(\mathcal{A}) = \mathcal{B}$.

For example, take $X = [0,1]$, $\mu = $ Lebesgue, $T(x) = 2x \bmod 1$ (the
doubling map), and $h(x) = x$.  The orbit $\{h \circ T^n\}(x) = \{2^n x \bmod 1\}$
generates $\mathcal{B}([0,1])$ as a $\sigma$-algebra (the dyadic structure separates
points).  But $\mathrm{span}\{2^n x \bmod 1 : n \geq 0\}$ is a proper subspace of
$L^2[0,1]$: it does not contain $x^2$.  So the linear span can be a proper dense
subspace even when $\sigma(\mathcal{A}) = \mathcal{B}$.

In the other direction, $h$ cyclic for $U_T$ (linear span dense) always implies
$\sigma(\mathcal{A}_h) = \mathcal{B}$ mod $\mu$ — see
Corollary~\ref{II:cor:cyclic-implies-reconstruction}.
\end{remark}


\begin{theorem}[Reconstruction theorem]
\label{II:thm:reconstruction}
Let $(X, \mathcal{B}, \mu, T)$ be an invertible measure-preserving system and
$h \in L^\infty(X, \mu)$.  The following are equivalent:
\begin{enumerate}[label=(\roman*)]
  \item $\mathcal{O}_h = \mathcal{B}$ modulo $\mu$;
  \item $\mathcal{A}_h$ is dense in $L^2(X, \mu)$;
  \item $\Phi_h : X \to \mathbb{R}^\mathbb{Z}$ is a measure-theoretic embedding:
    $\Phi_h^{-1}(\mathcal{B}(\mathbb{R}^\mathbb{Z})) = \mathcal{B}$ modulo $\mu$.
\end{enumerate}
\end{theorem}

\begin{proof}
\emph{(i) $\Leftrightarrow$ (ii).}
Apply the density bridge (Lemma~\ref{II:lem:density-bridge}) with
$\mathcal{A} = \mathcal{A}_h$ and $\sigma(\mathcal{A}_h) = \mathcal{O}_h$:
$\overline{\mathcal{A}_h}^{L^2} = L^2(X, \mu)$ iff $\mathcal{O}_h = \mathcal{B}$
mod $\mu$.

\emph{(i) $\Leftrightarrow$ (iii).}
By Definition~\ref{II:def:delay-map},
$\mathcal{O}_h = \Phi_h^{-1}(\mathcal{B}(\mathbb{R}^\mathbb{Z}))$.
So $\mathcal{O}_h = \mathcal{B}$ mod $\mu$ iff
$\Phi_h^{-1}(\mathcal{B}(\mathbb{R}^\mathbb{Z})) = \mathcal{B}$ mod $\mu$,
which is exactly condition~(iii).
\end{proof}

Condition~(iii) means $\Phi_h$ induces an isomorphism of measure algebras
$(X/\mu, \mathcal{B}/\mu) \cong (\mathbb{R}^\mathbb{Z}/(\Phi_h)_*\mu,
\mathcal{B}(\mathbb{R}^\mathbb{Z})/(\Phi_h)_*\mu)$.
Since $(\Phi_h(Tx))_n = h(T^{n+1}x) = (\sigma(\Phi_h(x)))_n$, the measure
$(\Phi_h)_*\mu$ is shift-invariant, and $(X, T, \mu)$ is isomorphic to
$(\mathbb{R}^\mathbb{Z}, \sigma, (\Phi_h)_*\mu)$ restricted to $\Phi_h(X)$.

Reconstruction is an asymptotic condition: the finite-lag algebras
$\mathcal{O}_h^{(L)} \uparrow \mathcal{O}_h$ are monotone, and the
separation defect $\delta(L) = \delta(\mathcal{O}_h^{(L)})$ measures
the gap at finite lag.  The rate at which $\delta(L)$ vanishes
determines the practical content of the theorem and depends on
mixing and geometric hypotheses beyond the scope of the present paper.

\begin{definition}[Cyclic vector]
\label{II:def:cyclic}
The function $h \in L^2(X, \mu)$ is a \emph{cyclic vector} for $U_T$ if
\[
  \mathcal{H}_h
  \;:=\;
  \overline{\mathrm{span}}\!\bigl\{U_T^n h : n \in \mathbb{Z}\bigr\}
  \;=\; L^2(X, \mu).
\]
\end{definition}

\begin{corollary}[Cyclic implies reconstruction]
\label{II:cor:cyclic-implies-reconstruction}
If $h \in L^\infty(X, \mu)$ is a cyclic vector for $U_T$, then
$\mathcal{O}_h = \mathcal{B}$ modulo $\mu$.
\end{corollary}

\begin{proof}
Since $\mathcal{A}_h$ contains $\mathrm{span}\{h \circ T^n\} =
\mathrm{span}\{U_T^n h\}$ as a subset, $\overline{\mathcal{A}_h}^{L^2}
\supseteq \mathcal{H}_h = L^2$.  By the density bridge, $\mathcal{O}_h =
\mathcal{B}$ mod $\mu$.
\end{proof}

\begin{remark}[The implication is strict]
\label{II:rem:strict}
The converse of Corollary~\ref{II:cor:cyclic-implies-reconstruction} fails in
general: $\mathcal{O}_h = \mathcal{B}$ does not imply $h$ is cyclic.
The doubling map example of Remark~\ref{II:rem:algebra-necessary} witnesses this.

The cyclic vector condition asks for density of the \emph{linear span} of the
orbit; reconstruction asks only for density of the \emph{algebra}.  For a bounded
$h$, the linear span is contained in the algebra, so cyclicity is the stronger
condition.

The two conditions coincide precisely when $U_T$ has \emph{simple
$L^2$-spectrum}: $L^2(X,\mu)$ is a cyclic $U_T$-module, meaning there exists
some cyclic vector.  In this case, every generating observation $h$ (satisfying
$\mathcal{O}_h = \mathcal{B}$) is automatically cyclic.  Simple spectrum is
a generic property of ergodic transformations in the Baire category sense
\cite{Halmos1944} but is not universal.
\end{remark}


The Stone space $\mathrm{St}(\mathcal{O}_h)$ of the observable algebra is the
compact Hausdorff space whose points are ultrafilters on $\mathcal{O}_h$, with
the natural Stone embedding $\iota : X \to \mathrm{St}(\mathcal{O}_h)$ sending
$x$ to the principal ultrafilter $\{E \in \mathcal{O}_h : x \in E\}$.

\begin{theorem}[Stone space identification]
\label{II:thm:stone-id}
Suppose $\mathcal{O}_h = \mathcal{B}$ modulo $\mu$.  Then $\iota$ is a
measure-theoretic isomorphism: $\iota$ identifies the measure algebra
$(X/\mu, \mathcal{B}/\mu)$ with
$(\mathrm{St}(\mathcal{O}_h)/\iota_*\mu,
\mathcal{B}(\mathrm{St}(\mathcal{O}_h))/\iota_*\mu)$.
\end{theorem}

\begin{proof}
\emph{Injectivity mod $\mu$.}  Suppose $\iota(x) = \iota(x')$.  For any
$E \in \mathcal{B}$ with $x \in E \not\ni x'$, find $F \in \mathcal{O}_h$
with $\mu(E \triangle F) = 0$.  Since $x \in F$ and $\iota(x) = \iota(x')$,
also $x' \in F$, so $x' \in (E \triangle F) \cup F^c$, placing $x'$ in a
null set.  Hence $\iota$ is injective $\mu$-a.e.

\emph{$\sigma$-algebra identification.}  The pullback $\iota^{-1}$ maps the
Baire $\sigma$-algebra of $\mathrm{St}(\mathcal{O}_h)$ (generated by the
clopen sets $[E] = \{u : E \in u\}$) to $\mathcal{O}_h = \mathcal{B}$ mod
$\mu$.

\emph{Measure recovery.}  Since $\mu$ is $\sigma$-additive on
$\mathcal{O}_h = \mathcal{B}$, the collective exhaustion condition
of~\cite{mahon_paper1} forces $\iota_*\mu$ to concentrate on the principal
ultrafilters $\iota(X)$, so $\iota_*\mu$ is supported on the image and the
measure algebras are identified.
\end{proof}

\noindent The Stone space~\cite{mahon_paper1} — built as a compact ambient
space for realising $\sigma$-additivity — is, under the reconstruction condition,
the state space $X$ itself.  The construction that began by seeking
a compact completion of the sample space ends by recognising it as the object
the observer was trying to recover.

\begin{example}[Cantor set: reconstruction without smoothness]
\label{II:ex:cantor}
Let $\mathcal{C} \subset [0,1]$ be the middle-thirds Cantor set, and let
$\mu_\mathcal{C}$ be the Cantor measure --- the pushforward of the fair-coin
product measure $\nu = (\tfrac{1}{2}\delta_0 + \tfrac{1}{2}\delta_2)^{\mathbb{N}}$
under the ternary coding map $\pi(d_1,d_2,\ldots) = \sum_{k \geq 1} d_k/3^k$.
Let $S : \mathcal{C} \to \mathcal{C}$ be the shift map
\[
  S(x) = \begin{cases} 3x & x \in \mathcal{C} \cap [0,\tfrac{1}{3}], \\
                        3x - 2 & x \in \mathcal{C} \cap [\tfrac{2}{3},1],
          \end{cases}
\]
and let $h(x) = x$.  The system $(\mathcal{C},\mu_\mathcal{C},S)$ is
measurably conjugate to the Bernoulli shift $(\{0,2\}^{\mathbb{N}},\nu,\sigma)$
via $\pi$; in particular $S$ is $\mu_\mathcal{C}$-invariant and ergodic.

For $x = \pi(d_1,d_2,\ldots)$ we have $S^n x = \pi(d_{n+1},d_{n+2},\ldots)$.
Since every point of $\mathcal{C}$ lies in $[0,\tfrac{1}{3}] \cup [\tfrac{2}{3},1]$,
\[
  d_{n+1} = 0 \iff S^n x \in \bigl[0,\tfrac{1}{3}\bigr], \qquad
  d_{n+1} = 2 \iff S^n x \in \bigl[\tfrac{2}{3},1\bigr],
\]
so $h \circ S^n$ recovers the $(n+1)$-st ternary digit of $x$.  The delay
vector $\Phi_h^{(L)}(x) = (x, Sx, \ldots, S^L x)$ reads off digits
$d_1,\ldots,d_{L+1}$; its fibres are the generation-$(L+1)$ cylinder sets
\[
  F_z = \{x \in \mathcal{C} : d_k(x) = z_k,\; 1 \leq k \leq L+1\},
  \quad z \in \{0,2\}^{L+1},
\]
each carrying $\mu_\mathcal{C}(F_z) = 2^{-(L+1)}$, so
$\mathcal{O}_h^{(L)} = \sigma(\Phi_h^{(L)}) = \sigma\{F_z : z \in \{0,2\}^{L+1}\}$.
Since $\mu_\mathcal{C}$-a.e.\ pair of distinct points differs at some finite
digit, these algebras separate points as $L \to \infty$:
\[
  \mathcal{O}_h \;=\; \sigma\!\Bigl(\bigcup_{L \geq 0} \mathcal{O}_h^{(L)}\Bigr)
  \;=\; \mathcal{B}(\mathcal{C}) \quad \text{mod } \mu_\mathcal{C}.
\]
By Theorem~\ref{II:thm:reconstruction}, $\Phi_h$ is a measure-theoretic embedding:
the finite-lag observable algebras refine to the full Borel structure as $L \to \infty$.

The observable algebra also encodes the geometry of $\mathcal{C}$ directly.
At each lag $L$ there are $2^{L+1}$ fibres of Euclidean diameter $(1/3)^{L+1}$,
so the box-counting ratio $\log 2^{L+1}/\log 3^{L+1} = \log 2/\log 3$ is an
algebraic invariant of the algebra itself, constant across all lags.  The Hausdorff
dimension is not a separate geometric calculation: it is read off from the refinement
rate of the observable algebra.

$\mathcal{C}$ is not a smooth manifold: it is compact and zero-dimensional,
with Hausdorff dimension $\log 2/\log 3 \approx 0.63$.  The differential-geometric
hypotheses of Takens's theorem --- $C^2$ ambient dynamics, integer manifold
dimension $d$, embedding bound $2d+1$ --- do not apply; nor do the finite-lag extensions of
Sauer--Yorke--Casdagli~\cite{SauerYorkeCasdagli1991}, which still require a
finite-dimensional ambient manifold or attractor, apply.
Theorem~\ref{II:thm:reconstruction} applies because it requires only that the
observable algebra be measurably generated; smoothness and finite dimensionality
play no role.
\end{example}

The classical Takens embedding theorem~\cite{Takens1981} reaches a related
conclusion by a different route: for a $C^2$ diffeomorphism on a compact
$d$-manifold and a generic smooth $h$, the finite-delay map
$\Phi_h^{(k)}(x) = (h(x), h(Tx), \ldots, h(T^{k-1}x))$ is a diffeomorphic
embedding into $\mathbb{R}^{2d+1}$ for $k \geq 2d+1$.
Sauer--Yorke--Casdagli~\cite{SauerYorkeCasdagli1991} extend this to
prevalence-generic observables and Lipschitz maps, while retaining the
finite-dimensional manifold hypothesis and the finite-lag, finite-dimension
embedding framework.  Botvinick-Greenhouse, Lunsford, and
Wang~\cite{BotvinickGreenhouse2025} develop a measure-theoretic
framework for delay embedding via optimal transport, replacing
topological with distributional equivalence.  The reconstruction
theorem here takes a different route: it requires only measurability,
uses all delays simultaneously, and replaces the dimension count with
the algebraic density condition; the embedding is measure-theoretic
rather than topological.  On a compact manifold
with generic smooth $h$, the orbit $\{h \circ T^n\}$ separates points, so
Stone--Weierstrass gives $\mathcal{A}_h$ uniformly dense in $C(M)$ and hence
$L^2$-dense: the density condition holds and the results are compatible.  For non-smooth or zero-dimensional systems — the Cantor example
being representative — the closer historical ancestor is the symbolic dynamics
and generating-partition tradition~\cite{Adler1998}: representing orbits
through successively finer finite partitions is exactly what the observable
algebra $\mathcal{O}_h^{(L)}$ does as $L$ grows, with faithfulness of $\Phi_h$
replacing the generating-partition condition.

\begin{remark}[Observable separation and Lyapunov exponents]
\label{II:rem:lyapunov}
When $T$ is $C^r$ and reconstruction holds, the delay map $\Phi_h$ carries a
natural instability rate: for $\mu$-a.e.\ $x$ and nearby $y$,
\[
  \lambda_h(x,y)
  \;:=\;
  \limsup_{n\to\infty} \frac{1}{n}
  \log d\!\bigl(\Phi_h(T^n x),\, \Phi_h(T^n y)\bigr).
\]
This \emph{observable separation exponent} measures how fast the reconstructed
dynamics distinguishes initial conditions through the observable, rather than
how fast tangent vectors grow.  When $h$ is generic in the sense of
Takens~\cite{Takens1981}, the finite-lag embedding $\Phi_h^{(L)}$ is
bi-Lipschitz on the attractor, sandwiching $\lambda_h$ between the top
Lyapunov exponent $\lambda_1$ from above and below; one expects
$\lambda_h = \lambda_1$ $\mu$-a.e., but making this precise requires
comparing the induced metric on $\Phi_h^{(L)}(X)$ with the ambient
Riemannian geometry as $L \to \infty$.
\end{remark}


% ====================================================================
\section{Discussion}
\label{II:sec:discussion}
% ====================================================================

The measure constructed in~\cite{mahon_paper1} already contains, without
further assumption, the structures developed here.  The Rokhlin
disintegration reads off a disintegration kernel for any pair
of observables; temporal indexing of that kernel forces
Chapman--Kolmogorov and hence the Markov semigroup; and in a
measure-preserving system the minimal sufficient factor becomes the
delay map.  Nothing is added to the measure --- only consequences are
drawn.

The mathematical content of §§\ref{II:sec:prediction}--\ref{II:sec:reconstruction} is standard once a $\sigma$-additive measure is in hand on the observable algebra; what the programme changes is the explanatory position of each object.  The individual components are classical.  Rokhlin disintegration~\cite{Rokhlin1952}
gives regular conditional distributions on standard Borel spaces; Halmos and
Savage~\cite{HalmosSavage1949} laid the measure-theoretic foundation for sufficient
statistics via the Radon--Nikodym theorem; Blackwell~\cite{Blackwell1953} placed
sufficiency in comparative terms; and Doob~\cite{Doob1953} established conditional
expectation, filtrations, and Markov semigroups as the language of modern probability.
Predictive-state representations~\cite{Littman2001} and observable operator
models~\cite{Jaeger2000} likewise distrust hidden-state parametrisations and encode
state through observable predictions.  Here the kernel, the factor, the semigroup,
and the delay map are not assumed, chosen, or built from a prior dynamical law ---
they are consequences of a single measure on an observable algebra, disclosed in
sequence.

The reconstruction theorem closes a circle opened in~\cite{mahon_paper1}.
The Stone space $\mathrm{St}(\mathcal{C})$ built there to realise
$\sigma$-additivity is an abstract completion of the observable algebra,
containing the realised state space as the principal ultrafilters but
not yet identified with it.  When $\Phi_h$ is faithful modulo~$\mu$
(Definition~\ref{II:def:faithful}), the three equivalent conditions of
Theorem~\ref{II:thm:reconstruction} jointly force $\sigma(\Phi_h) =
\mathcal{B}$ modulo~$\mu$, and the Stone space collapses onto $X$
itself: the abstract completion is the concrete state space.  The
question left open in~\cite{mahon_paper1} --- whether prediction,
dynamics, and state recovery follow from observation --- is answered
affirmatively, under faithfulness.  What faithfulness means in practice,
and whether it can be certified from finite data, remain open.

% ====================================================================
\ifstandalone
\bibliographystyle{plain}
\bibliography{references}
\fi
